\documentclass[12pt,a4paper,oneside]{article}
\usepackage[brazil]{babel}
\usepackage[utf8]{inputenc}
\usepackage{olymp}
\usepackage{color}
\usepackage[dvipsnames]{xcolor}
\usepackage{listings}

\setcounter{problem}{3}

\begin{document}

\begin{problem}{Elemento repetido}{repete}{2}
 
Dado um conjunto de listas, sua tarefa é deve determinar se há um elemento
repetido dentro de cada lista. Caso haja, indicar também qual é este elemento.

%\Notes
%
%Observações, se tiver.

\InputFile

A primeira linha da entrada contém um inteiro $T$ representando o total de listas que serão informadas. A seguir, virão $T$ linhas. Cada linha
inicia com um inteiro $N$ indicando o tamanho da lista informada nesta linha. Em seguida, na mesma linha, virão os $N$ valores $V_i$ da lista, todos inteiros.

Restrições: $1 \leq T, N \leq 1000$; $-10^6 \leq V_i \leq 10^6$.

\OutputFile

Seu programa deve gerar $T$ linhas de saída, uma para cada lista da entrada, contendo: `` \texttt{NAO}'', caso não haja elemento repetido na lista; ou ``\texttt{SIM:} $X$'', caso $X$ seja o elemento repetido. 

É garantido que, se houver repetição, apenas um elemento se repete.

%\Notes
%Nesta questão, seu programa deve usar a função \texttt{fatorial} dada %acima exatamente como está, não a modifique! Sua
%tarefa é apenas criar o programa com a função \texttt{main}.

\Example

\begin{example}
\exmp{
5
3 -90 54 7
5 34 67 34 -9 76
7 978 23 465 -12 543 -12 1
9 0 54 12 65 897 23 1 99 5
10 9 8 23 64 2 67 64 12 64 98
}{
NAO
SIM: 34
SIM: -12
NAO
SIM: 64
}%
\end{example}

%Comentários sobre o exemplo, se necessário.

\end{problem}

\end{document}
